\documentclass[11pt]{article}

\usepackage[a4paper,margin=1in]{geometry}
\usepackage[T1]{fontenc}
\usepackage[utf8]{inputenc}
\usepackage{lmodern}
\usepackage{amsmath,amssymb}
\usepackage{booktabs}
\usepackage{longtable}
\usepackage{array}
\usepackage{enumitem}
\usepackage{hyperref}
\usepackage{xcolor}

\hypersetup{
  colorlinks=true,
  linkcolor=blue!50!black,
  urlcolor=blue!50!black,
  citecolor=blue!50!black
}

\newcommand{\file}[1]{\texttt{#1}}
\newcommand{\ec}[1]{\texttt{#1}}
\newcommand{\todoec}{\textcolor{red!60!black}{\texttt{admit.}}}

\title{Proof Artifact Report for the HAETAE NTT EasyCrypt Development}
\author{Artifact documentation for the \file{proof/} directory}
\date{\today}

\begin{document}
\maketitle

\begin{abstract}
This report documents the EasyCrypt proof artifact currently stored in the
\file{proof/} directory of the HAETAE NTT repository. The artifact aims to
connect a Jasmin-extracted implementation of HAETAE polynomial arithmetic to
coefficient-level specifications over the finite field $\mathbb{Z}_{64513}$.
The development follows the proof architecture used in the ML-KEM/Kyber
artifacts, but it is adapted to HAETAE's modulus, 32-bit Montgomery arithmetic,
and 256-coefficient polynomial layout. The report explains the role of each
proof module, the main definitions and theorems, the current proof status, and
the operational checks available to test the artifact.
\end{abstract}

\section{Artifact Scope}

The \file{proof/} directory contains six substantive EasyCrypt modules and three
small array-wrapper files:
\begin{itemize}[leftmargin=1.5em]
\item \file{HAETAE\_poly\_extract.ec}: extracted EasyCrypt model of the Jasmin
      code, including the concrete twiddle tables and the imperative procedures
      for Montgomery reduction, field multiplication, NTT, inverse NTT, and
      coefficient-wise base multiplication.
\item \file{HAETAE\_GFq.ec}: finite-field model for HAETAE's modulus
      $q = 64513$.
\item \file{HAETAE\_Montgomery.ec}: generic signed Montgomery reduction theory,
      instantiated later with the HAETAE parameters.
\item \file{HAETAE\_FCLib.ec}: low-level word lemmas needed to bridge Jasmin
      bit-vector operations with mathematical integer semantics.
\item \file{HAETAE\_Fq.ec}: instantiation of Montgomery arithmetic for the
      extracted \ec{\_\_montgomery\_reduce} and \ec{\_\_fqmul} routines.
\item \file{HAETAE\_poly.ec}: abstract NTT specification, losslessness proofs,
      representation bridges, relational correctness statements, and wrapper
      correctness theorems.
\end{itemize}

The auxiliary files \file{Array256.ec}, \file{BArray1024.ec}, and
\file{WArray1024.ec} provide array interfaces used by the extracted code and
its proofs. The corresponding \file{.eco} files are compiled EasyCrypt
artifacts that speed up replay and serve as evidence that individual modules
have been checked in the local toolchain.

\section{Design Overview}

The proof development is layered. Each layer eliminates a different gap between
the extracted implementation and the target mathematical statement.

\subsection{Layer 1: Extracted Program Model}

\file{HAETAE\_poly\_extract.ec} defines \ec{module M}, which is the concrete
program artifact to be verified. Its key procedures are:
\begin{itemize}[leftmargin=1.5em]
\item \ec{M.\_\_montgomery\_reduce}: reduces a signed 64-bit intermediate to a
      32-bit coefficient using HAETAE's Montgomery parameters.
\item \ec{M.\_\_fqmul}: multiplies two signed 32-bit coefficients and applies
      Montgomery reduction.
\item \ec{M.\_poly\_ntt}: iterative breadth-first NTT on 256 coefficients using
      concrete twiddle tables \ec{jzetas}.
\item \ec{M.\_poly\_invntt}: inverse NTT using \ec{jzetas\_inv} followed by
      final scaling.
\item \ec{M.\_poly\_basemul}: coefficient-wise multiplication in the
      Montgomery domain.
\item Wrapper procedures \ec{poly\_ntt\_jazz}, \ec{poly\_invntt\_jazz}, and
      \ec{poly\_basemul\_jazz}.
\end{itemize}

This file is intentionally operational: it contains the concrete control flow,
loop counters, word-level loads/stores, and precomputed constants that later
proof layers must interpret.

\subsection{Layer 2: Algebraic Model}

\file{HAETAE\_GFq.ec} defines the coefficient type by cloning EasyCrypt's
\ec{ZModField} at modulus $64513$. The file also introduces the signed
projection \ec{as\_sint}, which is used throughout the artifact to express
range bounds on coefficients.

\file{HAETAE\_Montgomery.ec} then develops a parameterized theory
\ec{SignedReductions}. The central object is the signed Montgomery reduction
operator
\[
\ec{SREDC}(a),
\]
whose main theorem \ec{SREDCp\_corr} proves both a range bound
\[
-q \le \ec{SREDC}(a) < q
\]
under the standard input precondition, and the expected modular congruence
\[
\ec{SREDC}(a) \equiv a \cdot R^{-1} \pmod q.
\]
This generic theorem is the algebraic backbone for the later correctness proofs
of HAETAE's concrete field multiplication.

\subsection{Layer 3: Word-Level Bridging}

\file{HAETAE\_FCLib.ec} contains the bit-level lemmas needed to connect the
Jasmin right-shift and sign-extension operators to integer division. The most
important theorem is \ec{SAR\_sem32}, which is intended to justify that
arithmetic right shift by 32 bits matches signed division by $2^{32}$. This is
precisely the semantic bridge required by the extracted Montgomery reduction
code.

\subsection{Layer 4: Field Arithmetic Correctness}

\file{HAETAE\_Fq.ec} instantiates \ec{SignedReductions} with the HAETAE
parameters:
\[
q = 64513,\qquad R = 2^{32},\qquad q^{-1} \bmod R = 940508161,\qquad
R^{-1} \bmod q = 50386.
\]
It then states the two key arithmetic correctness lemmas:
\begin{itemize}[leftmargin=1.5em]
\item \ec{montgomery\_reduce\_corr\_h}: correctness of
      \ec{M.\_\_montgomery\_reduce}.
\item \ec{fqmul\_corr\_h}: correctness of \ec{M.\_\_fqmul}.
\end{itemize}
These theorems are lifted into probability-1 judgments
(\ec{phoare}) by combining them with losslessness results.

\subsection{Layer 5: Polynomial-Level Correctness}

\file{HAETAE\_poly.ec} is the main algorithmic proof file. It introduces:
\begin{itemize}[leftmargin=1.5em]
\item an abstract coefficient-level \ec{NTT} module with procedures
      \ec{ntt}, \ec{invntt}, and \ec{basemul};
\item losslessness theorems for the abstract algorithms and the extracted
      implementations;
\item representation bridges such as \ec{lift\_array256},
      \ec{zetas\_coeff}, \ec{zetas\_inv\_coeff}, and
      \ec{signed\_bound\_cxq};
\item algebraic bridge lemmas \ec{rrinvcoeff} and \ec{mul\_congr};
\item relational correctness statements
      \ec{poly\_ntt\_equiv}, \ec{poly\_invntt\_equiv}, and
      \ec{poly\_basemul\_equiv};
\item wrapper correctness statements for the public entry points.
\end{itemize}

This file is where the ML-KEM proof structure is being ported to HAETAE. The
current \ec{poly\_ntt\_equiv} statement already carries a stage-sensitive bound
invariant: a signed-$q$ input bound is tracked to a signed-$9q$ output bound
through the eight NTT stages.

\section{Module-by-Module Inventory}

\begin{longtable}{>{\raggedright\arraybackslash}p{0.22\textwidth}
                  >{\raggedright\arraybackslash}p{0.18\textwidth}
                  >{\raggedright\arraybackslash}p{0.42\textwidth}
                  >{\raggedright\arraybackslash}p{0.12\textwidth}}
\toprule
Module & Role & Key objects and theorems & Current status \\
\midrule
\endfirsthead
\toprule
Module & Role & Key objects and theorems & Current status \\
\midrule
\endhead
\file{HAETAE\_GFq.ec} &
Finite-field interface &
\ec{q}, \ec{prime\_q}, clone of \ec{ZModField}, signed projection
\ec{as\_sint} &
Complete in current file \\[0.4em]

\file{HAETAE\_Montgomery.ec} &
Generic signed Montgomery theory &
\ec{SignedReductions}, \ec{smod}, \ec{SREDC}, \ec{SREDCp\_corr} &
Core theorem complete in current file \\[0.4em]

\file{HAETAE\_FCLib.ec} &
Word-level helper lemmas &
\ec{mask32\_hi}, \ec{aux32\_0}, \ec{aux32\_1}, \ec{aux32\_2},
\ec{SAR\_sem32} &
One admitted lemma \\[0.4em]

\file{HAETAE\_Fq.ec} &
Concrete arithmetic correctness &
\ec{montgomery\_reduce\_corr\_h}, \ec{montgomery\_reduce\_corr},
\ec{fqmul\_corr\_h}, \ec{fqmul\_corr} &
Two admitted lemmas \\[0.4em]

\file{HAETAE\_poly\_extract.ec} &
Extracted imperative model &
\ec{module M}, \ec{jzetas}, \ec{jzetas\_inv}, \ec{\_poly\_ntt},
\ec{\_poly\_invntt}, \ec{\_poly\_basemul} &
Operational model present \\[0.4em]

\file{HAETAE\_poly.ec} &
Abstract spec and relational correctness &
\ec{module NTT}, \ec{lift\_array256}, \ec{signed\_bound\_cxq},
\ec{poly\_ntt\_equiv}, \ec{poly\_invntt\_equiv},
\ec{poly\_basemul\_equiv} &
Nine admitted lemmas \\
\bottomrule
\end{longtable}

\section{Main Mathematical Content}

\subsection{Finite-Field and Representation Bridge}

The coefficient type is represented abstractly as elements of
$\mathbb{Z}_{64513}$ and concretely as signed 32-bit integers stored in
\ec{BArray1024}. The bridge
\[
\ec{lift\_array256}(p) =
\ec{Array256.init}(\lambda i.\,\ec{incoeff}(\ec{W32.to\_sint}(\ec{get32}\ p\ i)))
\]
maps a concrete polynomial buffer to the abstract coefficient vector consumed by
the specification. This is the canonical relation used in all equivalence
theorems.

The twiddle tables require one additional adjustment. Since the extracted code
multiplies through \ec{\_\_fqmul}, the concrete table entries are interpreted in
the proof as
\[
\ec{incoeff}(\text{table entry}) \cdot R^{-1},
\]
captured by \ec{zetas\_coeff} and \ec{zetas\_inv\_coeff}. The lemma
\ec{rrinvcoeff} is then used to cancel the extra Montgomery factor.

\subsection{Losslessness}

Before relational correctness is attempted, the artifact establishes
losslessness of both the abstract algorithms and their extracted counterparts.
This matters for two reasons:
\begin{itemize}[leftmargin=1.5em]
\item it guarantees that the imperative loops are structurally well-founded in
      EasyCrypt's semantics;
\item it provides the probabilistic side conditions needed to lift Hoare-style
      arithmetic results into probability-1 correctness statements.
\end{itemize}

The files \file{HAETAE\_poly.ec} and \file{HAETAE\_Fq.ec} already contain
completed losslessness lemmas such as \ec{ntt\_ll}, \ec{poly\_ntt\_ll},
\ec{montgomery\_reduce\_ll}, and \ec{fqmul\_ll}.

\subsection{Montgomery Arithmetic}

The intended proof story for field arithmetic is standard:
\begin{enumerate}[leftmargin=1.5em]
\item prove that the extracted right shift and truncation operations implement
      the integer divisions and modular reductions appearing in the definition
      of \ec{SREDC};
\item instantiate the generic theorem \ec{SREDCp\_corr};
\item conclude that \ec{\_\_fqmul} computes multiplication in the Montgomery
      domain with the expected $R^{-1}$ factor.
\end{enumerate}

This bridge is already structurally in place, but its last steps still depend
on the admitted word-level and arithmetic lemmas in \file{HAETAE\_FCLib.ec} and
\file{HAETAE\_Fq.ec}.

\subsection{NTT-Level Correctness}

The central algorithmic theorem is currently:
\begin{quote}
\begin{verbatim}
equiv poly_ntt_equiv :
  NTT.ntt ~ M._poly_ntt :
    r{1} = lift_array256 rp{2} /\
    zetas{1} = zetas_coeff /\
    signed_bound_cxq rp{2} 0 256 1
    ==>
    res{1} = lift_array256 res{2} /\
    signed_bound_cxq res{2} 0 256 9.
\end{verbatim}
\end{quote}

This statement is important for paper exposition because it shows the exact
contract carried by the ported ML-KEM invariant:
\begin{itemize}[leftmargin=1.5em]
\item relational equality between the abstract NTT state and the lifted
      concrete state;
\item stage-sensitive signed bounds on the concrete representation;
\item explicit alignment between the extracted twiddle table and the abstract
      coefficient table.
\end{itemize}

The proof skeleton in \file{HAETAE\_poly.ec} already identifies the intended
HAETAE analogues of the ML-KEM ingredients:
\begin{itemize}[leftmargin=1.5em]
\item \ec{signed\_bound\_cxq} for concrete coefficient range tracking;
\item \ec{zeta\_bound} for twiddle-table range facts;
\item \ec{rrinvcoeff} to eliminate the residual Montgomery factor;
\item \ec{mul\_congr} to convert concrete integer multiplication into abstract
      coefficient multiplication.
\end{itemize}

\section{Current Proof Gaps}

At the time of writing, the artifact still contains twelve explicit
\todoec{} placeholders:
\begin{itemize}[leftmargin=1.5em]
\item \file{HAETAE\_FCLib.ec}: \ec{SAR\_sem32}.
\item \file{HAETAE\_Fq.ec}: \ec{montgomery\_reduce\_corr\_h},
      \ec{fqmul\_corr\_h}.
\item \file{HAETAE\_poly.ec}: \ec{zeta\_bound}, \ec{zetainv\_bound},
      \ec{ntt\_inner\_loop\_correct}, \ec{poly\_ntt\_equiv},
      \ec{poly\_ntt\_corr}, \ec{poly\_invntt\_equiv},
      \ec{poly\_invntt\_corr}, \ec{poly\_basemul\_equiv},
      \ec{poly\_basemul\_corr}.
\end{itemize}

These gaps are not cosmetic; they identify the precise frontier between the
already-ported framework and the remaining mechanization work.

\subsection{What Is Already Established}

Several important components are already in place:
\begin{itemize}[leftmargin=1.5em]
\item the finite-field model and parameter instantiation;
\item the generic signed Montgomery correctness theorem;
\item the extracted imperative model and all concrete constants;
\item losslessness proofs for the abstract and concrete procedures;
\item the shape of the main relational theorems and the HAETAE-specific NTT
      bound invariant.
\end{itemize}

\subsection{What Still Blocks Full End-to-End Correctness}

The remaining work clusters into three proof-critical areas:
\begin{enumerate}[leftmargin=1.5em]
\item \textbf{Word semantics.} The negative-case proof of \ec{SAR\_sem32} is
      still missing.
\item \textbf{Field arithmetic.} The high-level correctness of
      \ec{\_\_montgomery\_reduce} and \ec{\_\_fqmul} is stated but not yet fully
      discharged.
\item \textbf{Algorithmic relational proofs.} The NTT, inverse NTT, and base
      multiplication equivalence proofs are scaffolded but still admitted.
\end{enumerate}

\section{Operational Validation}

The repository now includes a dedicated artifact test runner:
\begin{center}
\file{tools/test-proof-artifacts.sh}
\end{center}
and matching \file{Makefile} targets:
\begin{itemize}[leftmargin=1.5em]
\item \ec{make test}
\item \ec{make proof-test}
\item \ec{make proof-test-quiet}
\item \ec{make proof-test-strict}
\end{itemize}

The test runner performs the following checks:
\begin{enumerate}[leftmargin=1.5em]
\item command availability (\ec{easycrypt}, \ec{why3}, \ec{why3server});
\item presence of required \file{.ec} and \file{.eco} artifacts;
\item freshness of compiled \file{.eco} files relative to their sources;
\item inventory of unresolved \todoec{} placeholders;
\item staged EasyCrypt smoke compiles for the proof modules.
\end{enumerate}

In the current local environment, the script reports an
\ec{ENVIRONMENT\_FAILURE} because \ec{why3server} is not available in
\ec{PATH}. This is an environment issue rather than a mathematical contradiction
in the proof artifact. The same run also reports that
\file{HAETAE\_poly.eco} is stale relative to \file{HAETAE\_poly.ec}, which is
expected after recent edits to the NTT invariant.

\section{How to Cite the Artifact State in a Paper}

For a paper or technical report, the artifact can be described accurately as
follows:
\begin{quote}
The EasyCrypt development contains a complete extracted model of HAETAE's
polynomial arithmetic, a finite-field formalization at $q = 64513$, a
parameterized Montgomery reduction theory, completed losslessness proofs for the
abstract and extracted procedures, and a partially mechanized port of the
ML-KEM relational proof architecture for NTT correctness. The remaining work is
localized to the word-level shift semantics, the final field-arithmetic bridge
lemmas, and the admitted relational correctness proofs for NTT, inverse NTT,
and base multiplication.
\end{quote}

This wording is precise about what has been formalized, what has been proved,
and what remains unfinished.

\section{Conclusion}

The \file{proof/} artifact is already organized as a serious verification
development rather than a flat collection of proof scripts. Its strongest
features are the clean separation between extracted code, algebraic theory, and
algorithmic correctness, together with the explicit reuse of the ML-KEM proof
structure. Its current limitations are also sharply localized and therefore
tractable. From a paper-writing perspective, this makes the artifact mature
enough to describe in detail, provided the write-up distinguishes clearly
between completed layers, admitted lemmas, and environment-dependent replay
requirements.

\end{document}
